\documentclass[11pt,a4paper]{article}

% =========================================================
% PACKAGES
% =========================================================
\usepackage[utf8]{inputenc}
\usepackage[T1]{fontenc}

\usepackage[margin=1in]{geometry}

\usepackage{amsmath,amssymb,amsfonts}
\usepackage{bm}
\usepackage{mathtools}

\usepackage{graphicx}
\usepackage{booktabs}
\usepackage{array}
\usepackage{multirow}
\usepackage{enumitem}

\usepackage{xcolor}
\usepackage{listings}

\usepackage{hyperref}
\usepackage[most]{tcolorbox}

\usepackage{titlesec}
\usepackage{setspace}
\usepackage{parskip}

\usepackage{caption}

% =========================================================
% HYPERREF
% =========================================================
\hypersetup{
    colorlinks = true,
    linkcolor  = blue,
    citecolor  = blue,
    urlcolor   = cyan
}

% =========================================================
% SPACING
% =========================================================
\setstretch{1.15}
\setlength{\parskip}{4pt}

% =========================================================
% LIST SETTINGS (global)
% =========================================================
\setlist{
    topsep     = 4pt,
    itemsep    = 2pt,
    parsep     = 0pt,
    leftmargin = 1.4em
}
\setlist[enumerate]{label=\arabic*.}
\setlist[itemize]{label=\textbullet}

% =========================================================
% SECTION FORMATTING
% =========================================================
\titleformat{\section}
    {\large\bfseries}
    {\thesection.}{0.6em}{}
\titlespacing{\section}{0pt}{14pt}{6pt}

\titleformat{\subsection}
    {\normalsize\bfseries}
    {\thesubsection}{0.6em}{}
\titlespacing{\subsection}{0pt}{10pt}{4pt}

\titleformat{\subsubsection}
    {\normalsize\itshape}
    {\thesubsubsection}{0.6em}{}
\titlespacing{\subsubsection}{0pt}{8pt}{3pt}

% =========================================================
% TABLE SETTINGS
% =========================================================
\renewcommand{\arraystretch}{1.35}
\setlength{\tabcolsep}{8pt}

% =========================================================
% TCOLORBOX STYLE
% =========================================================
\tcbset{
    defbox/.style={
        colback  = gray!6,
        colframe = black!55,
        arc      = 3pt,
        boxrule  = 0.6pt,
        left     = 8pt,
        right    = 8pt,
        top      = 6pt,
        bottom   = 6pt,
    }
}

% =========================================================
% CUSTOM COMMANDS
% =========================================================
\newcommand{\vect}[1]{\bm{#1}}
\newcommand{\R}{\mathbb{R}}

% =========================================================
% TITLE
% =========================================================
\title{%
    \vspace{-1.5cm}
    \textbf{\Huge Neuron}\\[0.4cm]
    \Large Mathematical Foundations of a Self-Organizing Cognitive Memory Graph
}

\author{%
    \textbf{Ajit Priyadarshi}\\
    \small IIT(BHU), Varanasi
}

\date{May 2026}

% =========================================================
% DOCUMENT
% =========================================================
\begin{document}

\maketitle
\vspace{-0.6cm}

% =========================================================
% ABSTRACT
% =========================================================
\begin{abstract}
Persistent, self-organizing memory is essential for autonomous AI agents to achieve long-term coherence. Traditional RAG systems often suffer from factual fragmentation and static retrieval. This paper introduces \textbf{Neuron}, a graph-based cognitive memory framework where beliefs (nodes) and associations (edges) evolve dynamically. We detail Neuron's multi-stage architecture—comprising surprise filtering, semantic abstraction, adaptive retrieval evolution, and a biologically-inspired ``Sleep Cycle'' for memory consolidation. We provide the mathematical foundations for novelty scoring, confidence reinforcement, and strategy optimization, offering a robust substrate for lifelong agent learning.
\end{abstract}

\vspace{0.4cm}

% =========================================================
% INTRODUCTION
% =========================================================
\section{Introduction}

Modern Large Language Models (LLMs) possess vast parametric knowledge but lack the
ability to autonomously update their worldview based on new experiences without
expensive fine-tuning. Traditional memory systems are typically implemented as flat
vector databases or short-term retrieval pipelines, resulting in:

\begin{itemize}
    \item Fragmented factual storage
    \item Weak contextual reasoning
    \item Redundant semantic accumulation
    \item Poor long-term coherence
\end{itemize}

Neuron addresses these limitations through a persistent graph-backed cognitive memory
architecture inspired by biological learning systems. Instead of treating memory as
isolated embeddings, Neuron organizes knowledge into a continuously evolving semantic
graph where:

\begin{itemize}
    \item \textbf{Nodes} represent distilled semantic beliefs
    \item \textbf{Edges} represent contextual associations
    \item \textbf{Retrieval} reinforces semantic pathways
    \item \textbf{Maintenance cycles} reorganize memory
\end{itemize}

The framework transforms memory from passive storage into a self-organizing cognitive
substrate capable of:

\begin{enumerate}
    \item Incremental learning
    \item Semantic abstraction
    \item Adaptive retrieval optimization
    \item Memory consolidation
    \item Long-term contextual reasoning
\end{enumerate}

By integrating vector similarity search with graph traversal dynamics and
reinforcement-based adaptation, Neuron provides a foundation for persistent autonomous
cognition.

% =========================================================
% DATA REPRESENTATION
% =========================================================
\section{Data Representation}

The fundamental unit of knowledge in Neuron is the \textbf{Belief}, represented as a
graph node. Semantic and contextual relationships are encoded as weighted graph edges.

% ---------------------------------------------------------
\subsection{Node Structure (Belief)}
% ---------------------------------------------------------

Each node $n$ is a structured semantic object containing the following attributes:

\begin{table}[h!]
\centering
\caption{Belief Node Structure}
\begin{tabular}{@{} p{4.2cm} p{8.8cm} @{}}
    \toprule
    \textbf{Field}         & \textbf{Description}                              \\
    \midrule
    ID                     & Unique UUID identifier                            \\
    User ID                & Multi-tenant ownership identifier                 \\
    Label                  & Distilled semantic fact                           \\
    Confidence $(C)$       & Certainty score in $[0,1]$                        \\
    Evidence Count         & Number of confirmations or observations           \\
    Contradiction Count    & Frequency of semantic conflicts identified        \\
    Domain Tags            & Semantic categories (e.g.\ AI, Finance)           \\
    Entities               & Extracted named entities                          \\
    Temporal Stability     & \{\texttt{stable}, \texttt{volatile},
                             \texttt{time-bound}\}                              \\
    Abstraction Level      & \{\texttt{specific}, \texttt{pattern},
                             \texttt{principle}\}                               \\
    Deprecated             & Boolean flag for resolution/forgetting            \\
    Embedding              & 384-dimensional semantic vector                   \\
    Created At             & Timestamp of node formation                       \\
    Last Confirmed At      & Timestamp of most recent evidence reinforcement   \\
    Last Contradicted At   & Timestamp of most recent conflict identified      \\
    \bottomrule
\end{tabular}
\end{table}

The semantic embedding is represented as:
\[
    E(n) \in \R^{384}
\]
using the default encoder \texttt{all-MiniLM-L6-v2}.

% ---------------------------------------------------------
\subsection{Edge Structure (Association)}
% ---------------------------------------------------------

Edges encode contextual relationships between beliefs. Each edge $e_{ij}$ contains:

\begin{table}[h!]
\centering
\caption{Association Edge Structure}
\begin{tabular}{@{} p{4.2cm} p{8.8cm} @{}}
    \toprule
    \textbf{Field}     & \textbf{Description}                                      \\
    \midrule
    ID                 & Unique edge UUID identifier                               \\
    From / To ID       & UUIDs of connected nodes                                  \\
    Relation Type      & \texttt{supports}, \texttt{contradicts},
                         \texttt{refines}, \texttt{depends\_on}, \ldots          \\
    Weight $(w)$       & Semantic strength in $[0.1,\,1.0]$                        \\
    Evidence Count     & Number of confirmations or traversals                     \\
    Created At         & Timestamp of edge formation                               \\
    Last Updated At    & Timestamp of most recent weight reinforcement             \\
    \bottomrule
\end{tabular}
\end{table}

% =========================================================
% SCORING TAXONOMY
% =========================================================
\section{Scoring Taxonomy and Input Mechanics}

Neuron's operation is governed by several adaptive scoring mechanisms.

\begin{table}[h!]
\centering
\caption{Neuron Scoring Metrics}
\label{tab:scores}
\begin{tabular}{@{} p{2.8cm} p{7.8cm} p{3.2cm} @{}}
    \toprule
    \textbf{Metric}           & \textbf{Formula / Logic}
                              & \textbf{Range / Threshold}          \\
    \midrule
    Novelty $(N)$
        & $\displaystyle 1 - \max\!\bigl(\operatorname{cosine\_sim}
            (V_{\!in},V_{\!neighbors})\bigr)$
        & $0.0$--$1.0$; novel if $N>0.3$          \\[8pt]

    Confidence $(C)$
        & $C_{t+1} = \min(0.95,\;C_t \times 1.1)$ per confirmation
        & $0.0$--$1.0$; valid if $C>0.35$         \\[8pt]

    Edge Weight $(w)$
        & $w_{t+1} = \min(1.0,\;w_t\times 1.05)$ per traversal
        & $0.1$--$1.0$; active if $w>0.6$         \\[8pt]

    Fitness $(F)$
        & EMA: $(1-\alpha)F_t + \alpha R$,\quad $\alpha=0.3$
        & $0.0$--$1.0$                             \\[8pt]

    Implicit Score $(S_i)$& $0.5 + 0.2(h>0) + 0.1(h>5)$
        & $0.5$--$0.8$                             \\
    \bottomrule
\end{tabular}
\end{table}

% ---------------------------------------------------------
\subsection{Novelty Score}
% ---------------------------------------------------------

Novelty quantifies semantic surprise relative to existing memory:
\[
    N \;=\; 1 - \max_{b \in B}\Bigl(\operatorname{cosine\_sim}(E(T),\,E(b))\Bigr)
\]
where $T$ is the incoming text, $E(\cdot)$ is the embedding function, and $B$ is the
current belief set.

% ---------------------------------------------------------
\subsection{Confidence Reinforcement}
% ---------------------------------------------------------

Confidence increases multiplicatively with repeated confirmations:
\[
    C_{t+1} \;=\; \min(0.95,\;C_t \times 1.1)
\]
This approach provides larger relative gains for low-confidence beliefs while decelerating as they approach the 0.95 saturation limit.

% ---------------------------------------------------------
\subsection{Fitness Score}
% ---------------------------------------------------------

Retrieval strategies are evaluated using an Exponential Moving Average (EMA):
\[
    F_{t+1} \;=\; 0.7\,F_t + 0.3\,R
\]
where $F_t$ is the prior strategy fitness and $R$ is the retrieval reward derived from:

\begin{enumerate}
    \item Explicit user feedback
    \item Implicit retrieval heuristics (see Section 3.4)
\end{enumerate}

% ---------------------------------------------------------
\subsection{Implicit Score}
% ---------------------------------------------------------

When explicit feedback is absent, retrieval quality is estimated via the implicit score $S_i$:
\[
    S_i \;=\; 0.5 + 0.2\,(h>0) + 0.1\,(h>5)
\]
where $h$ denotes the number of relevant belief nodes (hits) successfully retrieved. This mechanism is heavily used during maintenance cycles to evolve strategies in the background.

% =========================================================
% INGESTION
% =========================================================
\section{Ingestion and Surprise Filtering}

Neuron uses a staged ingestion pipeline designed to minimize redundancy while
preserving semantically valuable information.

% ---------------------------------------------------------
\subsection{Surprise Filter Gate}
% ---------------------------------------------------------

For every incoming text chunk $T$, the system computes:
\[
    S_{\max}
    \;=\;
    \max_{b \;\in\; \mathrm{NN}(T,\,5)}
    \Bigl(\operatorname{cosine\_sim}(E(T),\,E(b))\Bigr)
\]
where $\mathrm{NN}(T,5)$ denotes the five nearest semantic neighbors.

The ingestion pipeline is divided into three semantic bands.

% ---------------------------------------------------------
\subsubsection{Novelty Band}
% ---------------------------------------------------------

\textit{Condition:}\quad $1 - S_{\max} \geq 0.3$

\begin{itemize}
    \item Input is considered semantically novel
    \item Sent to the LLM abstraction pipeline
    \item A new belief node is created
\end{itemize}

% ---------------------------------------------------------
\subsubsection{Confirmation Band}
% ---------------------------------------------------------

\textit{Condition:}\quad $0.75 \leq S_{\max} \leq 0.90$

\begin{itemize}
    \item Existing node evidence count incremented
    \item Confidence reinforced
\end{itemize}

% ---------------------------------------------------------
\subsubsection{Redundancy Band}
% ---------------------------------------------------------

\textit{Condition:}\quad $S_{\max} > 0.90$

\begin{itemize}
    \item Input discarded
    \item Prevents semantic duplication
    \item Reduces memory bloat
\end{itemize}

% =========================================================
% DATA ADDITION METHODS
% =========================================================
\section{Data Addition Methods}

% ---------------------------------------------------------
\subsection{Single-Chunk Pipeline}
% ---------------------------------------------------------

The primary real-time ingestion pipeline consists of:

\begin{enumerate}
    \item \textbf{Surprise Check.}
          If $N < 0.3$, the system updates or skips existing beliefs.

    \item \textbf{LLM Abstraction.}
          Novel content is distilled into:
          \begin{itemize}
              \item Semantic labels
              \item Domain tags
              \item Entities
              \item Abstraction level
              \item Temporal stability
          \end{itemize}

    \item \textbf{Node Creation.}
          A new belief node is inserted into the graph.

    \item \textbf{Local Linking.}
          Edges are created to the two nearest neighbors, with
          $w = \operatorname{similarity}$.
\end{enumerate}

% ---------------------------------------------------------
\subsection{Fast Batch Mode}
% ---------------------------------------------------------

Designed for high-throughput ingestion, the system first computes the
self-similarity matrix:
\[
    M = V V^{\top}
\]
Any chunk $i$ satisfying $\exists\, j < i : M_{ij} > 0.95$ is discarded. Remaining
chunks are connected through $k$-NN graph construction ($k = 3$), with edge weights
initialized using cosine similarity.

% ---------------------------------------------------------
\subsection{Deep Batch Mode}
% ---------------------------------------------------------

Deep Batch mode prioritizes semantic abstraction and long-range linkage:

\begin{enumerate}
    \item \textbf{Windowed LLM Abstraction.}
          Default window size: $n = 20$.

    \item \textbf{Global Entity Linking.}
          For each extracted entity:
          \begin{enumerate}
              \item Search the graph for the strongest matching entity node
              \item Select the node with highest evidence count
              \item Create a cross-batch semantic edge with $w = 0.9$
          \end{enumerate}
\end{enumerate}

This creates persistent semantic hubs across datasets and time.

% =========================================================
% RETRIEVAL
% =========================================================
\section{Retrieval and Myelination}

Neuron retrieval combines vector similarity with graph traversal.

% ---------------------------------------------------------
\subsection{Graph Traversal}
% ---------------------------------------------------------

Retrieval proceeds in two stages:

\begin{enumerate}
    \item Vector seed selection ($k = 5$ seeds)
    \item Breadth-First graph traversal to depth $D = 3$
\end{enumerate}

Edges satisfying $w < 0.6$ are treated as high-resistance pathways and ignored.

% ---------------------------------------------------------
\subsection{Myelination}
% ---------------------------------------------------------

Traversed edges are reinforced after successful retrieval:
\[
    w_{t+1} \;=\; \min(1.0,\;w_t \times 1.05)
\]
This creates emergent semantic routing optimization analogous to biological
myelination.

% =========================================================
% ADAPTIVE EVOLUTION
% =========================================================
\section{Adaptive Strategy Evolution}

Neuron autonomously evolves retrieval strategies.

% ---------------------------------------------------------
\subsection{Selection and Exploration}
% ---------------------------------------------------------

The framework uses an $\epsilon$-greedy policy with $\epsilon = 0.2$:

\begin{itemize}
    \item \textbf{Explore:} random strategy selected with probability $\epsilon$
    \item \textbf{Exploit:} best-fitness strategy selected with probability $1 - \epsilon$
\end{itemize}

% ---------------------------------------------------------
\subsection{Evolution Cycle}
% ---------------------------------------------------------

Evolution is triggered every 50 retrieval events:

\begin{enumerate}
    \item Rank all strategies by fitness $F$
    \item Remove the lowest-fitness strategy
    \item Clone the best-performing strategy
    \item Apply stochastic mutation (see below)
\end{enumerate}

% ---------------------------------------------------------
\subsection{Mutation Functions}
% ---------------------------------------------------------

Child parameters are derived from the parent via:

\begin{align}
    k_{\text{child}}      &= \operatorname{round}\!\bigl(
                               k_{\text{parent}} + \mathcal{U}(-1,1)\bigr) \\[4pt]
    D_{\text{child}}      &= \operatorname{round}\!\bigl(
                               D_{\text{parent}} + \mathcal{U}(-1,1)\bigr) \\[4pt]
    w_{\min,\text{child}} &= w_{\min,\text{parent}} + \mathcal{N}(0,0.1)   \\[4pt]
    \tau_{\text{child}}   &= \tau_{\text{parent}}   + \mathcal{N}(0,0.1)
\end{align}

where $k$ is the retrieval seed count, $D$ the traversal depth, $w_{\min}$ the
minimum traversable edge weight, and $\tau$ the similarity threshold.

% =========================================================
% MAINTENANCE
% =========================================================
\section{Maintenance: The Sleep Cycle}

Neuron periodically enters a maintenance phase inspired by biological sleep
consolidation.

% ---------------------------------------------------------
\subsection{Definitions}
% ---------------------------------------------------------

\begin{tcolorbox}[defbox]
    \textbf{Maintenance} — Low-level graph operations:
    \begin{itemize}
        \item Pruning
        \item Consolidation
        \item Coherence repair
    \end{itemize}

    \medskip

    \textbf{Sleep Cycle} — High-level orchestration phase responsible for:
    \begin{enumerate}
        \item Implicit scoring
        \item Strategy evolution
        \item Graph maintenance
    \end{enumerate}
\end{tcolorbox}

Maintenance cycles execute every $24$ hours.

% ---------------------------------------------------------
\subsection{Internal Workflow}
% ---------------------------------------------------------

\subsubsection{1.\; Implicit Scoring}

Unscored retrieval events are evaluated using:
\[
    S_i \;=\; 0.5 + 0.2\,(\text{hits}>0) + 0.1\,(\text{hits}>5)
\]
The resulting scores update retrieval strategy fitness.

\subsubsection{2.\; Pruning}

Beliefs satisfying both of the following conditions are deprecated or removed:

\begin{itemize}
    \item No confirmation within the past 30 days
    \item $C < 0.35$
\end{itemize}

This minimizes graph entropy and preserves semantic efficiency.

\subsubsection{3.\; Consolidation}

When $\geq 5$ nodes share a common entity hub, the system synthesizes a higher-level
abstraction called a \textbf{Master Belief}. The Master Belief inherits:

\begin{itemize}
    \item The maximum confidence among constituents
    \item The sum of constituent evidence counts
\end{itemize}

This enables the organic emergence of generalized semantic concepts.

% =========================================================
% EMERGENT BEHAVIOR
% =========================================================
\section{Emergent Cognitive Behavior}

Through repeated retrieval, reinforcement, pruning, and consolidation, Neuron exhibits
several emergent properties:

\begin{enumerate}
    \item \textbf{Semantic Path Optimization.}
          Frequently traversed reasoning pathways become dominant.

    \item \textbf{Contextual Memory Clustering.}
          Related beliefs self-organize into semantic hubs.

    \item \textbf{Long-Term Concept Formation.}
          Repeated observations evolve into generalized principles.

    \item \textbf{Adaptive Retrieval Evolution.}
          Retrieval policies improve autonomously over time.

    \item \textbf{Dynamic Knowledge Compression.}
          Consolidation reduces redundancy while preserving semantic richness.
\end{enumerate}

% =========================================================
% CONCLUSION
% =========================================================
\section{Conclusion}

Neuron transforms memory from passive storage into a continuously evolving cognitive
substrate. By integrating:

\begin{itemize}
    \item Vector similarity search
    \item Graph-structured semantic relationships
    \item Reinforcement-driven pathway optimization
    \item Adaptive retrieval evolution
    \item Biologically-inspired consolidation mechanisms
\end{itemize}

the framework enables persistent self-organizing agent intelligence. Unlike traditional
RAG systems that retrieve isolated fragments, Neuron maintains an evolving semantic
worldview in which beliefs, associations, confidence, and abstraction continuously
co-adapt over time.

The architecture provides a foundation for future autonomous systems capable of:

\begin{enumerate}
    \item Lifelong learning
    \item Dynamic memory consolidation
    \item Contextual reasoning
    \item Adaptive strategy evolution
    \item Coherent long-term cognition
\end{enumerate}

As AI agents increasingly operate across extended time horizons and complex
environments, systems such as Neuron may become foundational infrastructure for
scalable machine cognition.

\vspace{0.8cm}
\begin{center}
    \textit{--- End of Paper ---}
\end{center}

\end{document}
